\documentclass[10pt,a4paper]{amsart}
\usepackage[margin=19mm]{geometry}
\usepackage{amsmath,amssymb,mathtools}
\usepackage[scr=rsfs,cal=euler]{mathalfa}
\usepackage{xfrac}
\usepackage[unicode,hidelinks,bookmarks=false]{hyperref}
\newcommand{\ie}{i.e.\ }
\newcommand{\cf}{cf.\ }
\newcommand{\resp}{respectively\ }
\newcommand{\Subsection}{\subsection}
\providecommand{\nobreakdash}{\nobreak\hbox{-}}
\setlength{\emergencystretch}{2em}
\multlinegap0pt
\usepackage{mdframed}
\usepackage{mdframed}
\usepackage{mdframed}
\usepackage{mdframed}
\usepackage{amssymb}
\usepackage{algorithm}\usepackage{algpseudocode}
\usepackage{tikz}
\usepackage{xcolor}
\usepackage{mdframed}
\usepackage{mathtools}
\usepackage[shortlabels]{enumitem}
\newtheorem{example}{Example}
\newtheorem{proposition}{Proposition}
\newtheorem{theorem}{Theorem}
\newcommand{\R}{\mathbb{R}}
\title{\bf On the Hidden Biases of Flow Matching Samplers}
\author{Soon Hoe Lim}
\date{Source: arXiv:2512.16768v3 (CC BY 4.0)}
\begin{document}
\maketitle

\noindent\emph{Source and licence.} \bf On the Hidden Biases of Flow Matching Samplers. Version arXiv:2512.16768v3 (CC BY 4.0), distributed by arXiv under the Creative Commons Attribution 4.0 International licence, http://creativecommons.org/licenses/by/4.0/, as recorded in the arXiv licence metadata for this exact version and shown on the abstract page https://arxiv.org/abs/2512.16768v3. Full licence notice: \texttt{licenses/arxiv-2512.16768-CC-BY-4.0.txt}.\par\smallskip

\noindent\textit{Editorial scope.} Counted unit: the article's proposition prop1 (source0.tex lines 523--537). The article's complete original proof of it is delivered with it (source0.tex lines 2062--2187, one \texttt{proof} environment of 126 lines, which the article places away from the statement). No mathematics is written, repaired or re-derived: the statement and the proof are the article's own lines, in the article's own notation, and no retained line is substituted or re-flowed.  Retained uncounted as the source-local context the proof needs: lines 440--443; lines 489--491; lines 503--518; lines 1015--1025; lines 1049--1075.  Nothing was omitted from any retained span, so the package carries no omission record.  No retained line contains a citation, so the wrapper carries no bibliography and no bibliography entry.  No retained line contains a figure environment, an image-inclusion command or drawing code, so no image is delivered and the package declares no asset dependency.  Editorial formatting only, in addition: an AMS \texttt{amsart} preamble that restates the article's own declarations and macros used by the retained lines, namely the environments \texttt{example}, \texttt{proposition}, \texttt{theorem} on separate counters, together with the standard packages those lines need.  The retained environments print in this document as Example 1, Proposition 1, Theorem 1 in order of appearance, whereas the article prints the same items with the numbering of its own full document.  The complete native source of the article as distributed in the arXiv e-print for this exact version is bundled unchanged as \texttt{sources/191274/source0.tex}.
\begin{align} \label{eq_ab}
    a_t(X) &= \frac{\frac{\partial \sigma_t}{\partial t}(X)}{\sigma_t(X)}, \quad 
    b_t(X) = \frac{\partial m_t}{\partial t}(X) - m_t(X) a_t(X).
\end{align}

\begin{equation} \label{eq_empRF}
    \hat{v}^*(t, z) = \sum_{i=1}^N w_i(t, z) \frac{x^{(i)}-z}{1-t},
\end{equation}

\begin{example}[Empirical Affine Flows and Smoothed Plug-in Targets]  \label{eg_empaffineflows}
The affine family also exhibits the smoothed plug-in regime in a particularly transparent way. Fix a  PDF $K > 0$ on $\R^d$, take $p_0(z)=K(z)$, and choose any $m_t$ and $\sigma_t$ such that
\[
m_0(X)=0, \qquad m_1(X)=X, \qquad \sigma_0(X)=1, \qquad \sigma_1(X)=\sigma_{\min}>0.
\]
Then the terminal conditional density is
\[
p_1(z\mid X=x^{(i)}) = \frac{1}{\sigma_{\min}^d}K\!\left(\frac{z-x^{(i)}}{\sigma_{\min}}\right),
\]
and averaging over the empirical target law yields the terminal marginal
\[
\tilde p_1(z)=\frac1N\sum_{i=1}^N p_1(z\mid X=x^{(i)})
=\frac{1}{N \sigma_{\min}^d} \sum_{i=1}^N K\!\left(\frac{z-x^{(i)}}{\sigma_{\min}}\right).
\]
Thus the terminal law is exactly the equally weighted kernel density estimator associated with kernel $K$ and bandwidth $\sigma_{\min}$. In particular, the affine-flow construction already contains a smoothed plug-in estimator of the target distribution. If $K$ is the standard Gaussian density, then this family converges formally to the rectified flow regime as the terminal bandwidth $\sigma_{\min} \downarrow 0$.
\end{example}

\begin{theorem}[Empirical setting, Gaussian source] \label{thm1}
Let $X_0 \sim \mathcal{N}(0, I_d)$ and suppose that we are given a fixed dataset $\mathcal{D}_N = \{x^{(i)}\}_{i \in [N]}$, $x^{(i)} \in \R^d$, with $M := \max_i \|x^{(i)}\| < \infty$.
Let $T \in [0,1)$ and define the instantaneous kinetic energy $K_t = \|\hat{v}^*(t, \psi_t(X_0))\|^2$ and the corresponding time-integrated kinetic energy $E_T = \int_0^T K_t \, dt$, where $\hat{v}^*$ is given in \eqref{eq_empRF} and $\psi_t$ solves $\dot{\psi}_t(X) = \hat{v}^*(t, \psi_t(X))$, $\psi_0(X) = X_0$, for $t \in [0, 1)$. Assume that there exists a unique solution to this ODE on $[0,T]$.

    \begin{itemize}
        \item[(a)] For each $t \in [0, T]$, there exist constants $C_t>0$, $c_t>0$, and threshold $U_t\ge 0$, depending only on $t$, $d$ and $M$, such that for every $u \ge U_t$,
        $$\mathbb{P}(K_t \ge u \mid \mathcal{D}_N) \le C_t \, e^{-c_t u}. $$
        \item[(b)] There exist constants $C_T>0$, $c_T>0$, and threshold  $U_T\ge 0$, depending only on $T$, $d$ and $M$, such that for every $u \ge U_T$,
        $$\mathbb{P}(E_T \ge u \mid \mathcal{D}_N) \le C_T \, e^{-c_T u}. $$
    \end{itemize} 
\end{theorem}

\begin{theorem}[Empirical setting, polynomial source-tail upper bound]
\label{thm2_polydecay}
Let $D_N = \{x^{(i)}\}_{i\in[N]}$ be a fixed  dataset with 
$\max_{i\in[N]} \|x^{(i)}\| < \infty$. Let $T \in [0,1)$.
Suppose the source distribution \(p_0\) satisfies the polynomial upper-tail bound:
\[
\mathbb{P}(\|X_0\| \ge s) \le \frac{C_\alpha}{s^\alpha}
\quad\text{for all } s \ge 1,
\]
for some constants $C_\alpha > 0$ and tail index $\alpha > 0$.

For the velocity field $\hat v^\ast$ defined in Proposition~\ref{prop1}, let
\[
A_{\max} := \sup_{t\in[0,T],\,i\in[N]} |a_t(x^{(i)})|,
\qquad
B_{\max} := \sup_{t\in[0,T],\,i\in[N]} \|b_t(x^{(i)})\|,
\]
and assume that \(A_{\max}<\infty\), \(B_{\max}<\infty\), and that there exists a unique solution to the ODE driven by $\hat v^\ast$
on $[0,T]$. Then, for each $t\in[0,T]$, there exist constants $C_t>0$ and threshold $U_t\ge 0$, depending only on $t, T, A_{\max}, B_{\max}, C_\alpha, \alpha$, such that for every $u\ge U_t$,
\[
\mathbb{P}(K_t \ge u \mid D_N) \le \frac{C_t}{u^{\alpha/2}}.
\]
Moreover, there exist constants $C_T>0$ and threshold $V_T\ge 0$, depending only on $T, A_{\max}, B_{\max}$, $C_\alpha$, $\alpha$, such that for every $u\ge V_T$,
\[
\mathbb{P}(E_T \ge u \mid D_N) \le \frac{C_T}{u^{\alpha/2}}.
\]
\end{theorem}

\begin{proposition} \label{prop1} 
For the family of affine flows in Example \ref{eg_empaffineflows}, the minimizer of the empirical FM objective over
\(L^2(dt\,\hat p_t(dz);\mathbb R^d)\) is unique \(dt\otimes \hat p_t\)-a.e. and, for a.e. \(t\), is given \(\hat p_t\)-a.e. by the closed-form formula:
    \begin{equation}
        \hat{v}^*(t, z) = \sum_{i=1}^N w_i(t,z) \cdot (a_t(x^{(i)}) z + b_t(x^{(i)})),
    \end{equation}
where $a_t$ and $b_t$ are given in \eqref{eq_ab}, and $w_i(t,z)$ is the kernel-dependent weighting function
\begin{equation}
    w_i(t,z) = \frac{p_t(z|x^{(i)})}{\sum_{j=1}^N p_t(z|x^{(j)})},
\end{equation}
with 
\begin{equation}
    p_t(z|x^{(i)}) = \frac{1}{\sigma_t^d(x^{(i)})} K\left(\frac{z-m_t(x^{(i)})}{\sigma_t(x^{(i)})} \right).
\end{equation}
\end{proposition}

\begin{proof}[Proof of Theorem~\ref{thm2_polydecay}]
The proof is analogous to that of Theorem~\ref{thm1}, with the Gaussian
tail bound replaced by the assumed power-law tail.

Recall from Proposition~\ref{prop1} that the empirical affine-flow minimizer
has the form
\[
\hat v^\ast(t,z)
= \sum_{i=1}^N w_i(t,z)\,\big(a_t(x^{(i)}) z + b_t(x^{(i)})\big),
\]
where the weights \(w_i(t,z)\) are nonnegative and sum to one. By the definition of
\[
A_{\max} := \sup_{t\in[0,T],\,i\in[N]} |a_t(x^{(i)})|,
\qquad
B_{\max} := \sup_{t\in[0,T],\,i\in[N]} \|b_t(x^{(i)})\|,
\]
we have, for all \(t \in [0,T]\) and all \(z \in \mathbb{R}^d\),
\begin{equation}
\label{eq:linear-growth-vhat}
\|\hat v^\ast(t,z)\|
\le \sum_{i=1}^N w_i(t,z)\big(|a_t(x^{(i)})|\,\|z\| + \|b_t(x^{(i)})\|\big)
\le A_{\max}\|z\| + B_{\max}.
\end{equation}

Let \(\psi_t\) denote the flow driven by \(\hat v^\ast\), i.e.,
\[
\dot\psi_t(X_0) = \hat v^\ast(t,\psi_t(X_0)), 
\qquad \psi_0(X_0) = X_0,
\]
and define \(r_t := \|\psi_t(X_0)\|\). Whenever\footnote{The same differential inequality holds for the upper Dini derivative of \(r_t\), which is sufficient for Grönwall.} \(r_t > 0\), we have, by the chain rule and
Cauchy--Schwarz inequality,
\[
\dot r_t
= \frac{\psi_t(X_0)}{\|\psi_t(X_0)\|} \cdot \dot\psi_t(X_0)
\le \|\hat v^\ast(t,\psi_t(X_0))\|.
\]
Using \eqref{eq:linear-growth-vhat} at \(z = \psi_t(X_0)\) gives
\[
\dot r_t
\le A_{\max} r_t + B_{\max}.
\]
By Grönwall’s lemma, there exist constants \(C_1(T), C_2(T) > 0\), depending only on
\(T, A_{\max}, B_{\max}\), such that for all \(t \in [0,T]\),
\begin{equation}
\label{eq:rt-bound}
r_t = \|\psi_t(X_0)\|
\le C_1(T)\,\|X_0\| + C_2(T).
\end{equation}

Define \(V_t := \hat v^\ast(t,\psi_t(X_0))\) and the instantaneous kinetic energy
\(K_t := \|V_t\|^2\). Combining \eqref{eq:linear-growth-vhat} and \eqref{eq:rt-bound},
we obtain
\[
\|V_t\|
\le A_{\max} r_t + B_{\max}
\le A_{\max}\big(C_1(T)\|X_0\| + C_2(T)\big) + B_{\max}
\le C_3(T)\,\|X_0\| + C_4(T),
\]
for suitable constants \(C_3(T), C_4(T) > 0\) depending only on \(T, A_{\max}, B_{\max}\).
Hence, by the inequality \((x+y)^2 \le 2(x^2 + y^2)\),
\begin{equation}
\label{eq:Kt-X0}
K_t
= \|V_t\|^2
\le 2C_3(T)^2\|X_0\|^2 + 2C_4(T)^2
\le C_K(T)\,\big(\|X_0\|^2 + 1\big),
\end{equation}
where we may take \(C_K(T) := 2\max\{C_3(T)^2, C_4(T)^2\}\).
Integrating \eqref{eq:Kt-X0} over \(t \in [0,T]\) yields the same type of bound
for the integrated kinetic energy
$E_T := \int_0^T K_t\,dt,$
i.e.,
\begin{equation}
\label{eq:ET-X0}
E_T
\le C_E(T)\,\big(\|X_0\|^2 + 1\big),
\end{equation}
for some constant \(C_E(T) := T C_K(T) > 0\) depending only on \(T, A_{\max}, B_{\max}\).

\medskip\noindent
\emph{Tail bounds.}
From \eqref{eq:Kt-X0}, for any \(u > 0\),
\[
\{K_t \ge u\}
\subseteq \Bigl\{\|X_0\|^2 \ge \frac{u}{C_K(T)} - 1\Bigr\}.
\]
Fix \(U_t\) large enough so that for all \(u \ge U_t\),
\(\frac{u}{C_K(T)} - 1 \ge 1\). Writing \(s := \sqrt{\frac{u}{C_K(T)} - 1}\),
we obtain
\[
\mathbb{P}\bigl(K_t \ge u \,\big|\, D_N\bigr)
\le \mathbb{P}\bigl(\|X_0\| \ge s \,\big|\, D_N\bigr)
= \mathbb{P}\bigl(\|X_0\| \ge s\bigr),
\]
since \(X_0\) is independent of \(D_N\). By the heavy-tailed assumption on \(p_0\),
for all \(s \ge 1\),
\[
\mathbb{P}\bigl(\|X_0\| \ge s\bigr)
\le \frac{C_\alpha}{s^\alpha}.
\]
For $u \ge U_t$ large enough so that $s^2 = \frac{u}{C_K(T)} - 1 \ge \frac{u}{2C_K(T)}$,
we have:
\[
\frac{1}{s^\alpha} \le \left(\frac{2C_K(T)}{u}\right)^{\alpha/2},
\]
and hence
\[
\mathbb{P}\bigl(K_t \ge u \,\big|\, D_N\bigr)
\le \frac{C_t}{u^{\alpha/2}},
\]
for all sufficiently large $u$, for a constant $C_t>0$ depending only on $t$, $T$, $A_{\max}$, $B_{\max}$, $\alpha$, and $C_\alpha$. This proves the first inequality in Theorem \ref{thm2_polydecay}.

The argument for $E_T$ is identical, using \eqref{eq:ET-X0} in place of
\eqref{eq:Kt-X0}. For any $u > 0$,
\[
\{E_T \ge u\}
\subseteq \Bigl\{\|X_0\|^2 \ge \frac{u}{C_E(T)} - 1\Bigr\},
\]
and the same substitution $s = \sqrt{\frac{u}{C_E(T)} - 1}$ together with the heavy-tailed bound on $\|X_0\|$ yields
\[
\mathbb{P}(E_T \ge u \,\big|\, D_N)
\le \frac{C_T}{u^{\alpha/2}},
\]
for all sufficiently large $u$, for a constant $C_T>0$ depending only on $T$, $A_{\max}$, $B_{\max}$, $\alpha$, and $C_\alpha$. This proves the second inequality in Theorem \ref{thm2_polydecay}.

\end{proof}

\end{document}
